%创建于2026-8-5 23:45
\documentclass{ctexart}
\usepackage[a4paper]{geometry}
\usepackage{mathtools}
\usepackage{amssymb}
\pagestyle{empty}
\begin{document}

\section*{数学}

\subsection*{一、初中（实数范围）}

\subsubsection*{1.相交线与平行线}
\begin{gather*}
\text{对顶角相等。} \\
\text{同位角相等，两条直线平行；两条直线平行，同位角相等。} \\
\text{内错角相等，两条直线平行；两条直线平行，内错角相等。} \\
\text{同旁内角互补，两条直线平行；两条直线平行，同旁内角互补。}
\end{gather*}
\subsubsection*{2.三角形}
\begin{gather*}
\text{三角形两边的和大于第三条边；三角形两边的差小于第三条边。} \\
\text{三角形有三条中线，这三条线交于一点，交点叫做三角形的重心。} \\
\text{三角形内角和等于$180^\circ$。} \\
\text{直角三角形的两个锐角互余；有两个角互余的三角形是直角三角形。} \\
\text{三角形的外角等于与它不相邻的两个内角的和。}
\end{gather*}
\paragraph{全等三角形}
\begin{gather*}
\text{全等三角形判定：SAS, ASA, AAS, SSS, HL。} \\
\text{角的平分线上的点到角的两边的距离相等；角的内部到角的两边距离相等的点在角的平分线上。}
\end{gather*}
\subsubsection*{3.轴对称}
\begin{gather*}
\text{成轴对称的两个图形中，连接对称点的线段被对称轴垂直平分。} \\
\text{经过线段中点并且垂直于这条线段的直线叫做垂直平分线。} \\
\text{线段垂直平分线上的点与这条线段两个端点的距离相等；} \\
\text{与线段两个端点距离相等的点在这条线段的垂直平分线上。} \\
\text{等腰三角形的两个底角相等；有两个角相等的三角形是等腰三角形。} \\
\text{等腰三角形底边上的中线、高及顶角平分线重合（简写成“三线合一”）。} \\
\text{等边三角形三个角都相等，并且每一个角都等于$60^\circ$。} \\
\text{三个角都相等的三角形是等边三角形。} \\
\text{有一个角是$60^\circ$的等腰三角形是等边三角形。} \\
\text{在直角三角形中，如果一个锐角等于$30^\circ$，那么它所对的直角边等于斜边的一半。}
\end{gather*}
\subsubsection*{4.幂运算}
\begin{gather*}
\text{设$a>0,b>0,m,n$为任意整数，且在$a^{\frac{m}{n}}=\sqrt[n]{a^m}$中$n>0$。}
\end{gather*}
\begin{align*}
a^0&=1 \\
a^{\frac{m}{n}}&=\sqrt[n]{a^m} \\
a^{-n}&=\frac{1}{a^n} \\
a^ma^n&=a^{m+n} \\
(ab)^n&=a^nb^n \\
(a^m)^n&=a^{mn} \\
\frac{a^m}{a^n}&=a^{m-n} \\
(\frac{a}{b})^n&=\frac{a^n}{b^n}
\end{align*}
\subsubsection*{5.因式分解}
\begin{align*}
(a\pm b)^2&=a^2\pm 2ab+b^2 \\
(a\pm b)^3&=a^3\pm 3a^2b+3ab^2\pm b^3 \\
a^2-b^2&=(a+b)(a-b) \\
a^2+b^2&=(a+b)^2-2ab=(a-b)^2+2ab \\
a^3-b^3&=(a-b)(a^2+ab+b^2) \\
a^3+b^3&=(a+b)(a^2-ab+b^2) \\
x^2+x(p+q)+pq&=(x+p)(x+q) \\
ax+ay+bx+by&=a(x+y)+b(x+y)=(a+b)(x+y)
\end{align*}
\subsubsection*{6.分式}
\begin{align*}
\frac{a}{b}\pm\frac{c}{b}&=\frac{a \pm c}{b}\quad(b\neq0) \\
\frac{a}{b}\pm\frac{c}{d}&=\frac{ad \pm bc}{bd}\quad(b\neq0,d\neq0) \\
\frac{a}{b}\cdot\frac{c}{d}&=\frac{ac}{bd}\quad(b\neq0,d\neq0) \\
\frac{\frac{a}{b}}{\frac{c}{d}}&=\frac{ad}{bc}\quad(b\neq0,d\neq0,c\neq0)
\end{align*}
\subsubsection*{7.二次根式}
\begin{gather*}
\text{在实数范围内，偶次根的被开方数必须为非负数。}
\end{gather*}
\begin{align*}
\sqrt{a}\sqrt{b}&=\sqrt{ab}\quad(a\ge0,b\ge0) \\
\sqrt{\frac{a}{b}}&=\frac{\sqrt{a}}{\sqrt{b}}=\frac{\sqrt{ab}}{b}\quad(a\ge0,b>0) \\
\sqrt{a^2}&=|a| \\
(\sqrt{a})^2&=a\quad(a\ge0)
\end{align*}
\subsubsection*{8.勾股定理}
\begin{gather*}
\text{如果直角三角形两条直角边长分别为$a, b$，斜边长为$c$，那么} \\
a^2+b^2=c^2\quad\sqrt{a^2+b^2}=c \\
c^2-b^2=a^2\quad\sqrt{c^2-b^2}=a \\
c^2-a^2=b^2\quad\sqrt{c^2-a^2}=b \\
\text{同理，满足以上条件的三角形是直角三角形。}
\end{gather*}
\paragraph{三角形的面积}
\begin{align*}
S&=\frac{ah}{2} \\
S&=\sqrt{p(p-a)(p-b)(p-c)}\quad p=\frac{a+b+c}{2} \\
S&=\sqrt{\frac{a^2b^2-(\frac{a^2+b^2-c^2}{2})^2}{4}}
\end{align*}
\subsubsection*{9.四边形}
\begin{gather*}
\text{四边形内角和等于$360^\circ$。} \\
\text{四边形外角和等于$360^\circ$。} \\
\text{$n$边形的内角和等于$(n-2)\times180^\circ$。} \\
\text{多边形的外角和等于$360^\circ$。} \\
\text{平行四边形的对边和对角相等。} \\
\text{平行四边形的对角线互相平分。} \\
\text{两组对边分别相等的四边形是平行四边形；} \\
\text{两组对角分别相等的四边形是平行四边形；} \\
\text{对角线互相平分的四边形是平行四边形。} \\
\text{一组对边平行且相等的四边形是平行四边形。} \\
\text{三角形的中位线平行于三角形的第三边，并且等于第三边的一半。} \\
\text{矩形就是长方形。} \\
\text{矩形的四个角都是直角；矩形的对角线相等。} \\
\text{直角三角形斜边上的中线等于斜边的一半。} \\
\text{对角线相等的平行四边形是矩形；有三个直角的四边形是矩形。} \\
\text{有一组邻边相等的平行四边形叫做菱形。} \\
\text{菱形的四条边都相等；} \\
\text{菱形的两条对角线互相垂直，并且每一条对角线平分一组对角。} \\
\text{对角线互相垂直的平行四边形是菱形；四条边相等的四边形是菱形。}
\end{gather*}
\subsubsection*{10.一元二次方程}
\begin{gather*}
\text{对于}ax^2+bx+c=0\quad(a\neq0) \\
\Delta=b^2-4ac \\
x=\frac{-b\pm\sqrt{\Delta}}{2a} \\
\end{gather*}
\begin{align*}
&\text{当$\Delta>0$时，方程有两个实数根。} \\
&\text{当$\Delta=0$时，方程有两个相等的根。} \\
&\text{当$\Delta<0$时，方程在实数范围无解。}
\end{align*}
\paragraph{韦达定理}
\begin{align*}
\text{对于}ax^2+bx+c&=0\quad(a\neq0) \\
x_1+x_2&=-\frac{b}{a} \\
x_1x_2&=\frac{c}{a}
\end{align*}
\subsection*{二、高中}
\subsubsection*{1.集合}
\begin{align*}
&\text{$\in$属于} \\
&\text{$\notin$不属于} \\
&\text{$\subseteq$子集（允许相等）} \\
&\text{$\subsetneq$真子集（不允许相等）} \\
&\text{$\nsubseteq$不包含于} \\
&\text{$\cup$并集} \\
&\text{$\cap$交集} \\
&\text{$\overline{A}$补集（需要指定全集）} \\
&\text{$\varnothing$空集} \\
&\text{$\mathbb{N}$自然数集} \\
&\text{$\mathbb{Z}$整数集} \\
&\text{$\mathbb{Q}$有理数集} \\
&\text{$\mathbb{R}$实数集} \\
&\text{$\mathbb{C}$复数集}
\end{align*}
\end{document}
